\section{Strict dimension for an actual proper closed integral image}
The maps here are the original scheme morphisms, and all whole dimensions
are computed from genuine original irreducible closed chains. No finite
dimension, maximal-chain or image-comparison certificate is an input.

\begin{lemma}[The whole original integral finite-type dimension is a natural number]
\label{lem:pl-gauss-original-whole-dimension-natural}
\uses{thm:pl-gauss-original-whole-dimension-equals-affine,scheme:actual-affine-finite-field-point}
\lean{EnvelopeProperty.PLGaussVerticalCurveSurfaceSelectionClosedDimension.exists_whole_dimension_eq_natCast}
\leanok
Let $k$ be any field and let the actual integral $Y$ have an actual
locally finite-type structure $p:Y\to\operatorname{Spec}k$. Then there
is a natural number $n$ with $\dim Y=n$ for the whole original scheme.
No global quasi-compactness or supplied finite-dimension bound is needed.
\end{lemma}
\begin{proof}
The actual integral scheme is nonempty. Choose an actual nonempty affine
neighbourhood of its genuine generic point. The original affine base map
makes its SAME actual coordinate ring a finite-type $k$-algebra. The
coordinate ring is nonzero because the affine is nonempty in an integral
scheme. Genuine Noether normalization gives a finite injective polynomial
algebra map into this ring. Incomparability bounds its prime-chain
dimension by the finite number of variables, while nonzeroness gives a
nonnegative dimension. Its actual ring dimension is therefore a natural
number, by the existing finite-type ring theorem. The proved whole
original scheme/affine dimension equality identifies that SAME dimension
with $\dim Y$. All chosen affine and algebra data are internally derived.
\end{proof}

\begin{theorem}[A nonsurjective actual closed integral image has strictly smaller whole dimension]
\label{thm:pl-gauss-original-proper-closed-image-strict-dimension}
\uses{lem:pl-gauss-original-whole-dimension-natural}
\lean{EnvelopeProperty.PLGaussVerticalCurveSurfaceSelectionClosedDimension.dimension_add_one_le_of_not_surjective}
\leanok
In the same original field-finite-type integral $Y$, let $D$ be integral
and let $c:D\hookrightarrow Y$ be an actual closed immersion. If the
actual point map of $c$ is not surjective, then
\[
 \dim D+1\le\dim Y.
\]
No finite dimension, chain, prime or prescribed image representative is
supplied. Local finite type of $D$ is derived from the actual composition.
\end{theorem}
\begin{proof}
The actual closed immersion is locally finite type; composing with the
original field structure gives the SAME $D$ a locally finite-type field
structure. By the preceding lemma $\dim D=n$ for an internally obtained
natural $n$. Krull dimension is the supremum of lengths of actual strict
chains of original irreducible closed subsets; the finite integer value
is attained. Choose such a genuine chain of length $n$.

The actual closed embedding sends an original irreducible closed subset
to its original closed image. This map on irreducible closeds is strictly
monotone, since the actual point map is inducing and injective. Thus the
whole chosen chain maps to a strict chain of the SAME length in $Y$.
Every image member lies in the actual closed range of $c$. The original
point map is nonsurjective, so this range is a proper subset of $Y$.
Consequently its last chain member is strictly below the whole original
$Y$, which is itself irreducible closed because $Y$ is integral. Append
this actual whole $Y$ to the mapped chain. This gives a genuine strict
chain of length $n+1$ in the original target, hence $n+1\le\dim Y$.
Substituting the actual original equality $\dim D=n$ proves the claim.
\end{proof}

\begin{corollary}[Equal successor dimensions force the actual closed map to be surjective]
\label{cor:pl-gauss-original-equal-successor-dimensions-surjective}
\uses{lem:pl-gauss-original-whole-dimension-natural,thm:pl-gauss-original-proper-closed-image-strict-dimension}
\lean{EnvelopeProperty.PLGaussVerticalCurveSurfaceSelectionClosedDimension.surjective_of_dimension_add_one_eq}
\leanok
For the same actual $p$ and original closed immersion $c:D\hookrightarrow Y$
with integral source and target, if their actual whole successor dimensions
satisfy $\dim D+1=\dim Y+1$, the actual point map $c$ is surjective.
This conditional dimension comparison will be discharged by the actual
cut-component dimension theorems in the horizontal-selection consumer.
\end{corollary}
\begin{proof}
If $c$ were nonsurjective, the preceding theorem would give
$\dim D+1\le\dim Y$. Substitute the given actual successor-dimension
equality. The first lemma gives $\dim Y=n$ for a genuine natural number,
so the resulting inequality is $n+1\le n$, impossible. This is actual
surjectivity of the original point map, not a replacement of its source,
image, kernel or scheme structure by a proxy.
\end{proof}

This strict-dimension theorem applies to an actual selected cut component
contained in an actual original central component. The two exact original
component-dimension theorems give equal successor dimensions. Surjectivity
then puts that central component's original generic-image point back in
the SAME cut, contradicting the cut's internally proved avoidance. Actual
horizontal selection, its coefficient dominance, normalization and the
surface induction are subsequent consumers, not conclusions here.
